\documentclass[../../main.tex]{subfiles}
\graphicspath{{../../image/}{../../figure/}}
\begin{document}

\section{注记}\label{section:ch15-注记}

用消失黏性法构造 $BV$ 解，长期以来一直是双曲守恒律系统理论中的一个核心公开问题。它最终由 Bianchini 与 Bressan [5] 以一种壮观的方式解决。本节的陈述是对那篇奠基性论文的节选。该理论的一份富有信息的综述见于 Bressan [14]。前期论文 Bianchini 与 Bressan [1,2,3,4] 为之奠定了基础。收敛速率见 Bressan 与 Yang [1]。该方法现已推广到双曲守恒律的初边值问题（Bianchini 与 Spinolo [1]，Spinolo [1]）。该方法的主导思想已卓有成效地用于构造 Riemann 问题解（Bianchini [6]），以及建立双曲守恒律半离散迎风格式（Bianchini [7,8]，Bressan、Baiti 与 Jenssen [1]）、对特殊系统的 Godunov 格式（Bressan 与 Jenssen [1]，Bressan、Jenssen 与 Baiti [1]）以及线性 Jin--Xin 松弛格式（Bianchini [9]）的收敛性。另见 Bianchini [4] 以及 Bressan 与 Shen [1]。同时，在 Godunov 格式与 Lax--Friedrichs 格式中也已检测到共振现象，它们可能驱使近似解的总变差趋于无穷；见 Bressan 与 Jenssen [1]，Baiti、Bressan 与 Jenssen [1]，以及 Bressan、Jenssen 与 Baiti [1]。今后还将有大量附加应用。还应强调：这些技术适用于一般拟线性严格双曲系统，无论其是否具守恒形式。当然，在非守恒情形，所构造的“解”未必在分布意义下满足方程，而应在非守恒激波理论（见\cref{section:ch15-注记}所引 LeFloch \emph{et al.}，并见 8.7 节）的语境中加以解释。

理解黏性激波如何形成与相互作用、以及黏性逼近如何收敛的最初步骤，由 Bressan 与 Donadello [1,2]、Shen 与 Park [1] 以及 Shen 与 Xu [1] 完成。

在 $L^1$ 中，解从\eqref{align:ch15-1513}到\eqref{align:ch15-1511}当黏性系数 $\mu$ 趋于零时的收敛速率，在 Bressan、Huang、Wang 与 Yang [1] 中讨论。

关于消失黏性方法的其他方面，文献浩繁。我们已在第六章看到该方法如何应用于 $L^\infty$ 或 $BV$ 框架下的标量守恒律。

一个重大公开问题是：当\eqref{align:ch15-1513}右端的单位矩阵被更一般的黏性矩阵替换时，该方法是否仍然适用。特别地，目前尚不知道是否可把（一维）等熵气体流的 Euler 方程实现为物理黏性系数消失的 Navier--Stokes 方程的解。

在消失黏性方法中，近似解 $U_\mu$ 携带关于黏性激波剖面的信息；当采用真实物理黏性时这一点尤其宝贵，但在 $\mu\to0$ 的极限中该信息丢失。这种信息丢失也发生在用消失毛细作用或松弛法构造解时，甚至发生在 8.7 节所勾勒的方法中——在该方法中，激波剖面本身决定弱解的概念。
\begin{enumerate}
\item\label{item:ch15-159-a} $\xi(\pm\infty,t)=\pm\infty$；
\item\label{item:ch15-159-b} $\xi(\cdot,t)$ 在 $U(\cdot,t)$ 的连续点集上可逆，且逆满足 $V(\xi^{-1}(x,t),t)=U(x,t)$；
\item\label{item:ch15-159-c} 若 $x$ 是 $U(\cdot,t)$ 的间断点，则 $s$ 在某闭区间，例如 $[s_-,s_+]$，上有 $\xi(s,t)=x$，且 $V(s_\pm,t)=U(x\pm,t)$，而 $V(\cdot,t)$ 在 $(s_-,s_+)$ 上追踪连接 $U(x-,t)$ 与 $U(x+,t)$ 的间断剖面。
\end{enumerate}

当 $(x,t)$ 是 $U$ 的近似跳跃间断点时，该间断剖面即为激波剖面；当 $(x,t)$ 属于不规则点集时，它即为完整波扇剖面。上述思想在概念上令人愉悦，并很可能找到技术应用。在该方向上见 Glass 与 LeFloch [1]。

\end{document}
